\documentclass[aspectratio=169]{beamer}

% ── Theme ────────────────────────────────────────────────────────────────────
\usetheme{default}
\usepackage{graphicx}
\usecolortheme{seagull}
\setbeamertemplate{navigation symbols}{}
\setbeamertemplate{frametitle}[default]
\setbeamercolor{frametitle}{fg=black}
\setbeamerfont{frametitle}{series=\bfseries, size=\normalsize}

% Global text margin (prevents content touching slide edge)
\setbeamersize{text margin left=7mm, text margin right=7mm}

% ── Colours ──────────────────────────────────────────────────────────────────
\definecolor{footlineblue}{RGB}{0,0,139}
\definecolor{lightblue}{RGB}{230,240,255}

% ── Block colours ────────────────────────────────────────────────────────────
\setbeamercolor{block title}{bg=footlineblue!80, fg=white}
\setbeamercolor{block body}{bg=lightblue, fg=black}
\setbeamerfont{block title}{size=\small, series=\bfseries}
\setbeamerfont{block body}{size=\small}

% ── Alert block ──────────────────────────────────────────────────────────────
\setbeamercolor{alerted text}{fg=footlineblue}

% ── Footline ─────────────────────────────────────────────────────────────────
\setbeamertemplate{footline}{%
  \leavevmode
  \hbox{%
    \begin{beamercolorbox}[wd=.40\paperwidth,ht=3.2ex,dp=1.0ex,leftskip=2mm]{footline}%
      \color{white}\small\insertsection
    \end{beamercolorbox}%
    \begin{beamercolorbox}[wd=.20\paperwidth,ht=3.2ex,dp=1.0ex,center]{footline}%
      \color{white}\small\insertframenumber/\inserttotalframenumber
    \end{beamercolorbox}%
    \begin{beamercolorbox}[wd=.40\paperwidth,ht=3.2ex,dp=1.0ex,rightskip=2mm,right]{footline}%
      \color{white}\small Min-Hop LEO Constellation Design
    \end{beamercolorbox}%
  }%
  \vskip0pt
}

% ── Packages ─────────────────────────────────────────────────────────────────
\usepackage{amsmath, amssymb, graphicx, booktabs, array}
\usepackage{hyperref}
\hypersetup{colorlinks=false, pdfborder={0 0 0}}

% ── Image placeholder ────────────────────────────────────────────────────────
% Usage: \imgbox{height}{caption}   (width always = \linewidth of its container)
\newcommand{\imgbox}[2]{%
  \fbox{%
    \parbox[c][#1][c]{\dimexpr\linewidth-2\fboxsep-2\fboxrule}{%
      \centering\scriptsize\color{gray}\textit{#2}%
    }%
  }%
}

% ── List spacing ─────────────────────────────────────────────────────────────
\setlength{\leftmargini}{1.2em}

% ── Title metadata ───────────────────────────────────────────────────────────
\title[Min-Hop Constellation Design]{%
  Minimum-hop Constellation Design\\[2pt]
  for Low Earth Orbit Satellite Networks}
\author[Karthik]{Karthik \quad {\small (22EC39017)}}
\institute[IIT Kharagpur]{%
  Department of Electronics and Electrical Communication\\
  IIT Kharagpur\\[2pt]
  {\small Supervisor: Prof.\ Rajarshi Roy}}
\date{Spring Semester, 2025--26}

\begin{document}

% ════════════════════════════════════════════════════════════════════════════
% 01 — Title
% ════════════════════════════════════════════════════════════════════════════
\begin{frame}
  \titlepage
\end{frame}

% ════════════════════════════════════════════════════════════════════════════
% 02 — Outline
% ════════════════════════════════════════════════════════════════════════════
\begin{frame}{Outline}
  \vspace{0.5cm}
  \small
  \begin{enumerate}
    \setlength\itemsep{0.55em}
    \item Motivation
    \item Problem Statement
    \item System Model
    \item Performance Metrics
    \item Topology Families
    \item Key Results
    \item Conclusions \& Future Work
  \end{enumerate}
\end{frame}

% ════════════════════════════════════════════════════════════════════════════
\section{Motivation}
% 03 — Motivation
% ════════════════════════════════════════════════════════════════════════════
\begin{frame}{Motivation}
  \vspace{0.1cm}
  \begin{columns}[T, totalwidth=\textwidth]

    \column{0.56\textwidth}
    \begin{block}{Context}
      \small
      \begin{itemize}
        \setlength\itemsep{0.35em}
        \item LEO mega-constellations use \textbf{steerable optical ISL terminals}.
        \item Link partner can be changed post-launch — topology is a \textbf{design choice}.
        \item Each extra hop adds delay, forwarding overhead, and ISL load.
      \end{itemize}
    \end{block}

    \vspace{0.25cm}

    \begin{block}{The Gap}
      \small
      \begin{itemize}
        \setlength\itemsep{0.35em}
        \item The \textbf{mesh grid} is used universally.
        \item It has \emph{never} been proved optimal.
      \end{itemize}
    \end{block}

    \column{0.40\textwidth}
    \vspace{0.1cm}
    \begin{figure}[htbp]
    \centering
    \includegraphics[width=1\textwidth,height=4cm]{LEO.png}
    \caption{LEO Satellite Network and Mesh Grid.}
    \label{fig:my_image}
\end{figure}

  \end{columns}
\end{frame}

% ════════════════════════════════════════════════════════════════════════════
\section{Problem Statement}
% 04 — Problem Statement
% ════════════════════════════════════════════════════════════════════════════
\begin{frame}{Problem Statement}
  \vspace{0.15cm}
  \small

  \begin{block}{Given}
    An $n_s \times n_o$ LEO constellation; each satellite maintains exactly $\Delta$ ISL connections.
  \end{block}

  \vspace{0.2cm}

  \begin{block}{Objective}
    Find the ISL topology that \textbf{minimises Average Shortest Path Length (ASPL) and Worst Shortest Path Length (WSPL)}.
  \end{block}

  \vspace{0.2cm}

  \begin{columns}[T, totalwidth=\textwidth]
    \column{0.47\textwidth}
    \begin{block}{Vertex-Symmetric}
      Same connection pattern at every node.\\[5pt]
      Simpler to manage.\\
      Bounded by $\Theta(\sqrt{N})$ ASPL.
    \end{block}

    \column{0.47\textwidth}
    \begin{block}{General Regular}
      Only degree constraint imposed.\\[5pt]
      Asymmetry exploited.\\
      Achieves $\Theta(\log N)$ ASPL.
    \end{block}
  \end{columns}
\end{frame}

% ════════════════════════════════════════════════════════════════════════════
\section{System Model}
% 05 — System Model
% ════════════════════════════════════════════════════════════════════════════
\begin{frame}{System Model}
  \vspace{0.1cm}
  \begin{columns}[T, totalwidth=\textwidth]

    \column{0.52\textwidth}
    \begin{block}{Constellation as a Torus}
      \small
      \begin{itemize}
        \setlength\itemsep{0.35em}
        \item $N = n_s \times n_o$ satellites on a lattice.
        \item Orbital wrap-around $\Rightarrow$ \textbf{unit-area torus}.
        \item ISL topology: $\Delta$-regular undirected graph $G = (V,E)$.
      \end{itemize}
    \end{block}

    \vspace{0.3cm}
    \small
    \textbf{Torus shift operator:}
    \[
      v \oplus e \;=\; \begin{pmatrix}
        v[0]+e[0] \bmod n_s \\[2pt]
        v[1]+e[1] \bmod n_o
      \end{pmatrix}
    \]

    \column{0.43\textwidth}
    \begin{figure}[htbp]
    \centering
    \includegraphics[width=0.8\textwidth]{TORUS.png}
    \caption{Torus representation of LEO Satellites.}
    \label{fig:my_image}
\end{figure}

    \vspace{0.25cm}

  \end{columns}
\end{frame}

% ════════════════════════════════════════════════════════════════════════════

% ════════════════════════════════════════════════════════════════════════════
\section{Topology Families}
% 08 — Offset Topology
% ════════════════════════════════════════════════════════════════════════════
\begin{frame}{The Offset Topology}
  \vspace{0.1cm}
  \begin{columns}[T, totalwidth=\textwidth]

    \column{0.52\textwidth}
    \begin{block}{One Change to the Mesh Grid}
      \small
      Replace cross-link $[0,1]^\top$ with $[\varepsilon,1]^\top$:
      \[
        \varepsilon \;=\; \left\lfloor \sqrt{\tfrac{2n_s}{n_o}} - 1 \right\rceil
      \]
      \begin{itemize}
        \setlength\itemsep{0.3em}
        \item Intra-orbital links kept intact.
        \item Cross-link distance $\propto 1/\!\sqrt{n_s}$ — \textbf{shortens as density grows}.
      \end{itemize}
    \end{block}

    \vspace{0.2cm}

    \begin{alertblock}{\small No Hardware Changes Required}
      \small Only the ISL terminal \textbf{pointing angle} changes.
    \end{alertblock}

    \column{0.43\textwidth}
    \begin{figure}[htbp]
    \centering
    \includegraphics[width=0.8\textwidth]{offsettorus.png}
    \caption{Torus with Offset.}
    \label{fig:my_image}
\end{figure}

    \vspace{0.2cm}
  \end{columns}
\end{frame}

% ════════════════════════════════════════════════════════════════════════════
% 09 — General Regular Topologies
% ════════════════════════════════════════════════════════════════════════════
\begin{frame}{General Regular Topologies}
  \vspace{0.1cm}
  \begin{columns}[T, totalwidth=\textwidth]

    % First Column: SA
    \column{0.48\textwidth}
    \begin{block}{Simulated Annealing (SA)}
      \small
      \begin{itemize}
        \setlength\itemsep{0.3em}
        \item \textbf{Init:} Steger--Wormald random regular graph.
        \item \textbf{Procedure:}
        \begin{enumerate}
            \item Propose a random double-edge swap.
            \item Evaluate the change in Average Shortest Path Length (ASPL).
            \item Accept swap if ASPL improves, or with probability $e^{-\Delta/T}$ if it worsens.
            \item Decrease temperature $T$ and repeat until convergence.
        \end{enumerate}
        \item \textbf{Result:} Achieves ASPL within \textbf{1.01--1.05$\times$} lower bound.
      \end{itemize}
    \end{block}

    % Second Column: PRGS
    \column{0.48\textwidth}
    \begin{block}{PRGS Algorithm}
      \small
      \begin{itemize}
        \setlength\itemsep{0.3em}
        \item \textbf{Init:} Torus space with constrained link range ($r \leq 0.25$).
        \item \textbf{Procedure:}
        \begin{enumerate}
            \item Partition the torus space into localized cells.
            \item Distribute vertices across these cells.
            \item Randomly generate intra-cell edges within the local groups.
            \item Form inter-cell edges while strictly maintaining the $r \leq 0.25$ range constraint.
        \end{enumerate}
        \item \textbf{Result:} Provably achieves $O(\log N)$ ASPL.
      \end{itemize}
    \end{block}

  \end{columns}
\end{frame}
% ════════════════════════════════════════════════════════════════════════════
\section{Results}
% 10 — Degree-4 Results
% ════════════════════════════════════════════════════════════════════════════
\begin{frame}{Results — Degree-4 Symmetric ($n_o = 4$)}
  \vspace{0.1cm}
  \begin{columns}[T, totalwidth=\textwidth]

    \column{0.50\textwidth}
    \begin{block}{ASPL and WSPL Comparison}
      \footnotesize
      \renewcommand{\arraystretch}{1.25}
      \begin{tabular}{@{}rr rr rr@{}}
        \toprule
        $n_s$ & $N$
              & \multicolumn{2}{c}{Mesh Grid}
              & \multicolumn{2}{c}{Sqrt-Offset} \\
        \cmidrule(lr){3-4}\cmidrule(lr){5-6}
              &     & ASPL & WSPL & ASPL & WSPL \\
        \midrule
        40  & 160 & 11.1 & 32 & 6.2  & 10 \\
        80  & 320 & 21.1 & 62 & 8.9  & 14 \\
        120 & 480 & 31.1 & 92 & 10.9 & 19 \\
        \bottomrule
      \end{tabular}
    \end{block}

    \vspace{0.2cm}

    \column{0.45\textwidth}
    \begin{figure}[htbp]
    \centering
    \includegraphics[width=1\textwidth,height=4cm]{ASPL.png}
    \caption{ASPL Comparision}
    \label{fig:my_image}
\end{figure}

  \end{columns}
\end{frame}

% ════════════════════════════════════════════════════════════════════════════
% 11 — WSPL Analysis
% ════════════════════════════════════════════════════════════════════════════
\begin{frame}{WSPL Analysis — A Novel Perspective}
  \vspace{0.1cm}
  \begin{columns}[T, totalwidth=\textwidth]

    \column{0.50\textwidth}
    \begin{block}{Summary Across All Families}
      \footnotesize
      \renewcommand{\arraystretch}{1.25}
      \begin{tabular}{@{}lrrr@{}}
        \toprule
        Topology & ASPL & WSPL & Ratio \\
        \midrule
        Mesh grid            & 31.1 & 92 & 13.6$\times$ \\
        Sqrt-offset          & 10.9 & 19 &  4.8$\times$ \\
        SA-refined           &  3.6 &  7 &  1.05$\times$ \\
        SA (dist.\ $r{=}0.25$) & 3.9 & 9 & 1.11$\times$ \\
        PRGS                 &  5.5 & 10 &  1.67$\times$ \\
        \bottomrule
      \end{tabular}
    \end{block}

    \vspace{0.2cm}

    \begin{block}{Key Insight}
      \small WSPL improvement is \textbf{systematically larger} than the ASPL improvement across all topology families.
    \end{block}

    \column{0.45\textwidth}
   \begin{figure}[htbp]
    \centering
    \includegraphics[width=1\textwidth, height=4cm]{WSPL_no.png}
    \caption{WSPL Comparision}
    \label{fig:my_image}
\end{figure}

  \end{columns}
\end{frame}

% ════════════════════════════════════════════════════════════════════════════
\section{Conclusion}
% 12 — Thank You
% ════════════════════════════════════════════════════════════════════════════
\begin{frame}
  \vfill
  \centering
  \Huge \textbf{Thank You!}
  \vfill
\end{frame}

\end{document}